\section{Background and Motivation}
\label{sec:background}

\subsection{Statute encoding in the wild}
\todo{Three threads to cover, each in 1--2 paragraphs.}

\paragraph{Logic-programming origins.} Sergot et al.~\cite{sergot1986british}
encoded the British Nationality Act as a Prolog program in 1986. The
encoding was complete enough to run case-style queries but exposed the
gap between informal drafting and Horn-clause semantics. Subsequent work
(Hammond, Bench-Capon) extended the approach without changing its shape.

\paragraph{Markup-driven codification.} \emph{Akoma Ntoso}~\cite{akomantoso}
and \emph{LegalRuleML}~\cite{legalruleml} took the orthogonal route of
heavyweight XML schemas designed for archival and inter-government exchange.
They encode \emph{structure} (who promulgated, when amended, which paragraph)
faithfully but have no native way to express element decomposition or
penalty arithmetic.

\paragraph{Modern DSLs.}
\begin{itemize}
\item \textbf{Catala}~\cite{catala2021}, designed by Merigoux et al. for
the French tax code, introduces \emph{prioritised default rules} as a
first-class construct, which we borrow for statutory exceptions
(\S\ref{subsec:exceptions}).
\item \textbf{lam4}~\cite{lam4} models contracts and regulations with a
deontic logic core and a notion of named-norm references / \texttt{IS\_INFRINGED}
predicate that we plan to adopt for inter-section reasoning.
\item \textbf{LexScript}~\cite{lexscript} (MIT) compiles statutes via a
tree-sitter parser and offers Tarjan-SCC analysis over reference graphs ---
a technique we adopt directly for cross-section reachability checks.
\item \textbf{LDOC}~\cite{ldoc} prioritises a Word-and-PDF authoring loop
and is the inspiration for our DOCX transpile target.
\end{itemize}

\subsection{Why a new DSL?}
\todo{Position Yuho in the design space. Summary table appears in
\S\ref{sec:related}; here, give the prose argument.}

Existing DSLs each cover a partial slice of what statute-encoding requires.
Catala has world-class default-rule semantics but no native penalty
combinators or illustration blocks. lam4 has a strong deontic core but
expects regulations, not penal sections. Akoma Ntoso has structure without
semantics. LexScript has clean tooling but a thin grammar. None target a
common-law penal code at full coverage.

Yuho takes the position that a \emph{statute-shaped} DSL --- where the
top-level construct is \texttt{statute N \{ definitions, elements,
penalty, illustrations, exceptions \}} --- is more honest than retrofitting
a contract or regulation DSL. A penal section has its own grammar; the
encoding should match it.

\subsection{Singapore Penal Code 1871: a tractable testbed}
\label{subsec:pc1871}
\todo{One paragraph: 524 sections, English-language drafting, public via
SSO, modest cross-reference density, illustration blocks ubiquitous,
Anglo-Indian heritage shared with India, Pakistan, Malaysia, Brunei.}
The Singapore Penal Code descends from Macaulay's 1860 Indian Penal Code
draft and retains its idiom: short numbered sections, inline illustrations,
penalty clauses with named combinators (``or both''), and Chapter IV
\emph{General Exceptions} that operate as cross-cutting defences. This
shape is what Yuho's grammar tracks.
